\chapter{Legal, Social, Ethical and Professional Issues}

\section{Legal Issues}

The legal dimension of this project concerns both the handling of research data during evaluation and the extent to which a consent-revocation interface can realistically support safer photo sharing in messaging environments. In the evaluation, the study used preset example imagery and scenario-based tasks rather than participants' own intimate content, and the prototype itself was implemented as a front-end simulation with in-memory state only, without real user accounts, server-side storage, or a persistent database. This meant the project was designed to minimise the collection and retention of personal data. In that sense, the approach is aligned with the data-minimisation and research-data handling principles associated with the UK GDPR and the Data Protection Act 2018, while avoiding any stronger claim that the prototype has been independently verified as legally compliant in the way a deployed commercial platform would need to be \cite{gdpr_2016, dpa_2018}.

At the design level, the UK GDPR is relevant not because it was written specifically for consent revocation in private messaging, but because it provides a useful benchmark for what meaningful consent and withdrawal should look like in digital systems. Under the GDPR, consent should be freely given, specific, informed and unambiguous, and withdrawal should be as easy as giving consent \cite{gdpr_2016}. This supports the aim of the Smart Photos prototype, which was designed to make permissions clearer before sharing and revocation more accessible after sharing. However, as Politou et al.\ note, revocation operates prospectively rather than retroactively, meaning that later withdrawal cannot guarantee complete erasure once data have already been circulated \cite{politou_2018}. This is especially relevant in WhatsApp-style photo sharing, where recipient-side saving remains possible even if the platform changes future in-app access. The project also relates to the Online Safety Act 2023, since it explores interface mechanisms intended to reduce harm and improve user control but it should still be understood as an academic prototype rather than a system that itself fulfils real platform duties \cite{online_safety_act_2023}. Finally, the prototype intentionally adopts familiar WhatsApp-like interaction patterns for evaluative realism. It remains as an original prototype and should not imply any affiliation with Meta.

\section{Social and Sustainability Issues}
Photo sharing in private messaging is a social process shaped by audience and changing interpersonal boundaries, not just a technical exchange. Petronio and Child \cite{petronio_cpm_2020} argue that privacy is negotiated through ongoing boundary management rather than fixed at a single moment, which is highly relevant to this project because consent may change after a photo has already been shared. The social scope of the problem is also wider than a single sender-recipient pair. Cherubini et al. \cite{cherubini_2021} show that photo sharing can create multi-party privacy conflicts when other depicted people are affected by a disclosure, and this becomes even more complex when a photo is forwarded to secondary recipients. In response, the Smart Photos prototype was designed to make these social relationships more visible by showing forwarding context, propagating permission changes across forwarded copies, and allowing depicted co-subjects to request revocation. These features could support more responsible photo-sharing norms by making the social responsibilities that concern consent revocation more explicit. However, the evaluation also showed that revocation can still be interpreted as mistrust or offence depending on the relationship and context. For this reason, the project offers one way to make communication practices in digital messaging safer and more accountable.

From a sustainability perspective, the environmental footprint of the project itself is minimal by design. The prototype was deployed on Vercel but operates entirely as a front-end application with in-memory state, with no persistent server, database, or cloud infrastructure of its own. Development and testing were conducted on a single personal machine using lightweight tools. At a systemic level, however, integrating a permission model of this kind into a production messaging platform would introduce additional computational overhead, since per-photo permission metadata and real-time permission updates would need to be stored, synchronised, and queried across devices at scale. Any platform considering adoption should therefore weigh the social benefits of clearer, revocable photo sharing against the energy cost of maintaining that state at production scale, as also noted in the discussion of platform providers in Section~\ref{sec:production-limitations}.

\section{Ethical Issues}
Researching consent revocation in private photo sharing involves a privacy-sensitive topic, since participants are asked to reflect on interpersonal boundaries and the potential misuse of shared images in messaging environments. Although the project does not focus on explicitly sexual content, engaging human participants still required careful ethical safeguards to protect their privacy and wellbeing. In line with the British Psychological Society's guidance for internet-mediated research, the evaluation was designed to minimise discomfort and avoid unnecessary disclosure of personal material \cite{bps_2021}. Participation was restricted to adults aged 18--24, informed consent was obtained at the start of each session, participants were told that they could withdraw at any time. To reduce ethical risk further, the study did not use participants' own photos or message histories. Instead, it relied on fictional screenshot-based scenarios, and a front-end prototype designed for evaluation. Interview recordings were transcribed for analysis, and participant data were anonymised in the reporting process. This approach is also consistent with value-sensitive design principles, which argue that systems should also be developed with attention to human values like privacy and autonomy \cite{friedman_hendry_2019}. Ethically, the prototype was therefore designed both to explore stronger user control and to avoid overstating what it can achieve after delivery.

\section{Professional Issues}

Throughout this project, the work was guided by the BCS Code of Conduct and Code of Good Practice, particularly its emphasis on acting in the public interest, maintaining professional competence and integrity, fulfilling relevant legal and ethical responsibilities, and contributing responsibly to the profession \cite{bcs_2022}. These principles are directly relevant to a project concerned with privacy-sensitive photo sharing, since poorly designed consent features could mislead users or overstate the level of protection available. In response, the Smart Photos model was developed as a research prototype grounded in literature and iterative design reasoning with findings reported cautiously within the scope of the project. This was especially important because the project deliberately focused on interaction design rather than real platform infrastructure, and therefore does not claim to solve technical issues.

Professional integrity in this project also required transparency and accountability in both implementation and reporting. The prototype was evaluated through a clear research process, supported by automated testing and explicit discussion of limitations, including the restricted scope of a front-end simulation. The report does not present revocation as complete control. It acknowledges that post-sharing permissions are partial and prospective, and it treats expectation management as a core design responsibility. This approach matches computing ethics arguments that professional responsibility means building useful systems and being honest about their limits and wider social consequences \cite{gotterbarn_2018}. 

The project also carries economic and commercial relevance. Clearer and more trustworthy consent features could strengthen user confidence in messaging platforms and improve their appeal in a competitive digital market. For platform providers, more responsible photo-sharing controls may also reduce reputational risk as privacy expectations grow. The project contributes to the profession by offering evidence-based design recommendations for platform providers, educators, and policymakers on how these new permission models might be handled more responsibly in messaging environments.